\providecommand{\Res}[2]{\qty[round-mode=figures,round-precision=3]{\fpeval{#1}}{#2}}
\providecommand{\Val}[1]{\num[round-mode=figures,round-precision=4]{\fpeval{#1}}}
\providecommand{\ResSci}[2]{\qty[round-mode=figures,round-precision=3,exponent-mode=scientific]{\fpeval{#1}}{#2}}
\subsubsection*{Question a) Largeur de la fente}
\begin{enumerate}[label=\arabic*)]
    \item \textbf{Variables connues et recherchées :}
    \begin{itemize}[label= -]
        \item $\lambda = \qty[exponent-mode=threshold]{\vsi{lam}}{m}$, $L = \qty{\v{L}}{m}$, $W = \qty[exponent-mode=threshold]{\vsi{W}}{m}$, donc $y_1 = W/2 = \Res{\vsi{W}/2}{m}$
        \item Variable recherchée : $a$
    \end{itemize}
    \item \textbf{Équations pertinentes :} pour de petits angles, $\sin\theta_1 \approx y_1/L$ et
    \[ a\sin\theta_1 = \lambda \]
    \item \textbf{Isolation de la variable recherchée :}
    \[ a = \frac{\lambda L}{y_1} = \frac{2\lambda L}{W} \]
    \item \textbf{Résolution numérique :}
    \[ a = \frac{2 \cdot \num[exponent-mode=threshold]{\vsi{lam}} \cdot \num{\v{L}}}{\num[exponent-mode=threshold]{\vsi{W}}} = \Res{1e6*2*\vsi{lam}*\v{L}/\vsi{W}}{\micro\metre} \]
    \item \textbf{Réponse :} la fente mesure \textbf{\Res{1e6*2*\vsi{lam}*\v{L}/\vsi{W}}{\micro\metre}} de large.
\end{enumerate}

\subsubsection*{Question b) Troisième minimum}
\[ \sin\theta_3 = \frac{3\lambda}{a} = \frac{3 \cdot \num[exponent-mode=threshold]{\vsi{lam}}}{\Val{2*\vsi{lam}*\v{L}/\vsi{W}}} \quad\Rightarrow\quad \theta_3 = \Res{asind(3*\vsi{W}/(2*\v{L}))}{\degree} \]
Le troisième minimum est à \textbf{\Res{asind(3*\vsi{W}/(2*\v{L}))}{\degree}}.

\subsubsection*{Question c) Nombre de minimums}
\begin{enumerate}[label=\arabic*)]
    \item \textbf{Variables connues et recherchées :} $a/\lambda = \Res{2*\v{L}/\vsi{W}}{}$ ; variable recherchée : $m_{\max}$
    \item \textbf{Équations pertinentes :} $\sin\theta_m = m\lambda/a \le 1$
    \item \textbf{Isolation de la variable recherchée :} $m \le a/\lambda$
    \item \textbf{Résolution numérique :} $m \le \Val{2*\v{L}/\vsi{W}}$, donc $m_{\max} = \num{\fpeval{floor(2*\v{L}/\vsi{W})}}$
    \item \textbf{Réponse :} on peut observer \textbf{\num{\fpeval{floor(2*\v{L}/\vsi{W})}}} minimums de chaque côté.
\end{enumerate}
